\documentclass{article}
\usepackage[utf8]{inputenc}
\usepackage{url}
\usepackage[spanish,english]{babel}
\usepackage{csquotes}
\usepackage{caption}
\usepackage[style=apa,backend=biber]{biblatex}
\addbibresource{references.bib}

\title{Openess to decentralization }
\author{Laura Itzel Rodríguez Dimayuga}
\date{November 2026}

\begin{document}

\maketitle

\section{Introduction}
Many Large Language Models (LLMs) are claimed to be open, but they are not. The openness of LLMs is often limited to the ability to access the model through the pretrain model, but the underlying arquitecture and its training data are not open, the are called instead \textit{open weight models}. Currently almost all local LLMs are just open weight models, which make it hard trace the origin of the model, its training data, and its limitations. 
\\ \\
This lack of transparency can lead to issues of accountability, bias, and misuse. It is necesary for the scientific community to be able to use, study, modify and audit the systems that mediate their knowledge, decisions and relationships. When we talk about democratazing technology, we are talking about the ability to have access to the means of production, and in this case, the means of production are the LLMs, not just use them. That is distribution of power, not redistribution of power, we need owenership of the means of cognitive production. 

\section{Problem}

\textcite{sapkota2025comprehensive} examine more than a hundred recent models and find that many of those advertised as open source do not report their training data or code, a practice that the literature calls \textit{open-washing}. OLMo \cite{groeneveld-etal-2024-olmo} and Pythia \cite{biderman2023pythia} show that full openness is feasible, because they publish data, training code, and intermediate checkpoints.



\section{Corpus}

\subsection{Conceptual Distinction}

Establish the definition of open weight, separete from full openness and propose four capabilities as criteria: use, study, modify, and audit. Weights guarantee the first and only partially the other three.

\subsection{Regulatory Framework}

Althoug \textcite{EU2024_1689} exempts certain transparency obligations for general-purpose models with free licenses whose weights, architecture, and usage information are public, it acknowledges that this release does not reveal substantial information about the training data. Therefore, it maintains the obligation to publish a summary of the content used and a copyright policy. 
\\ \\
\textcite{qiu2026modeling} build on this ambiguity and model how the openness threshold set by the regulator alters developers' publication decisions.
Show that the European Regulation defines openness by the published artifacts and not by the capabilities it enables, and that it recognizes this limitation by retaining obligations on data. 
\\ \\
\textcite{manchanda2026open} argue, from the capabilities approach, that open development enhances transparency, reproducibility, community participation, and auditability. They propose that an AI is democratic when it redistributes power, both in its development and in its action.

\subsection{Philosophical and political axes}
\begin{itemize}
\item Use the Marxist theory of dependency to analyze the problem from a Latin American perspective \textcite{reyes2025pattern}.
\item If social interaction is mediated by LLMs, it cannot be deliberated about what cannot be inspected. \textcite{manchanda2026open} treat LLMs as infrastructure that mediates access to knowledge, and \textcite{machado2025toward} associate open models with independent scrutiny and bias mitigation.
\item Previous work describes power as the ability to modulate attention through shared scripts. LLMs infer attributes of their users and condition their behavior on those inferences, including refusals to respond, and this internal state can be read and modified. This case nuances the thesis: audit is already possible with access to weights. What weights do not provide is the provenance of the data or the possibility of remaking the model.
\item Research shows that the personality exhibited by an LLM can be measured and shaped. \textcite{maeda2024human} analyze how anthropomorphic design induces trust and parasocial bonds. Whoever controls the model also controls the tone of the bond, and users of a closed system cannot verify for what purposes it is adjusted.
\item \textcite{habermas2022neuer} argue that democracy lives in public debate, but what happens when the public sphere has conditions that prevent interaction with users? The public sphere is a space of communication where citizens can debate and deliberate on matters of common interest. The emergence of bots and their interaction with users can affect the quality of debate and deliberation. Bots can be used to spread false or biased information, which can affect public opinion and decision-making, increasing hate speech, polarization, and disinformation.
\end{itemize}
\subsection{Objections} 

\begin{itemize}
\item Opening data and processes increases the risk of misuse.
\item Fully open models still perform worse
\item Opening the data does not decentralize anything if computing remains concentrated.
\item The complexity of managing and maintaining open-source projects can be a significant burden.
\end{itemize}



\begin{table}[t]
\centering
\caption{Comparison of LLMs from a Marxist perspective. The table presents four key concepts from Marxist theory and their corresponding implications for LLMs, along with questions that arise from these implications.}
\label{tab:llm_comparison}
\renewcommand{\arraystretch}{1.2}
\begin{tabular}{|p{0.18\textwidth}|p{0.36\textwidth}|p{0.36\textwidth}|}
\hline
Term & LLMs & Questions \\
\hline
Means of production & Data, GPUs, code, models, cloud infrastructure & Who owns the resources to train and deploy AI? \\
\hline
Surplus value / extraction & Personal data, annotation labor, attention and engagement, cognitive capital & What value is extracted from interactions with the models? \\
\hline
Private  property & Weights under restrictive licenses, closed datasets, APIs & Is “openness” real, or is it a form of marketing? \\
\hline
Alienation & Dependence on opaque systems for thinking, writing, or making decisions & Does the user understand and control the tool they use? \\
\hline
\end{tabular}
\end{table}





%bibliography
\printbibliography
\end{document}